\chapter{Bitácora de decisiones}
\label{cap:decisiones}

Las decisiones que un lector nuevo cuestionaría, con la alternativa que se
descartó y por qué. La mayoría están argumentadas en los mensajes de commit y en
los comentarios del propio código, que es de donde salen.

\begin{pktable}[caption={Decisiones de arquitectura},label={tab:adr}]{colspec={X[0.55,l]X[2,l]X[1.6,l]X[2.6,l]}}
  ID & Decisión & Alternativa descartada & Razón \\
  \pkid{AD-01} & Dos esquemas en una base, migrados por herramientas distintas. & Un solo esquema y un solo migrador. & Cada esquema lo define el sistema que sabe cómo debe ser. Better Auth crea la tabla que un plugin nuevo necesita; Alembic no puede saberlo. El riesgo (que autogenerate proponga borrar lo ajeno) se cierra con dos filtros. \\
  \pkid{AD-02} & Carta y especie son capas separadas, unidas por el \pkcode{dexId} del origen. & Una tabla con el nombre del Pokémon dentro de la carta. & Una carta es una impresión y una especie es lo retratado; unirlas por nombre falla con formas regionales, cartas de entrenador y traducciones. \\
  \pkid{AD-03} & El escaneo puntúa todas las señales en lugar de filtrar por alguna. & Filtrar por número de colector o por set. & Un campo mal leído descartaría la carta correcta. Puntuando, un HP erróneo cuesta 0.05 y la respuesta sigue apareciendo. \\
  \pkid{AD-04} & \pkcode{CommittingRoute} confirma antes de que la respuesta salga. & Confiar en el commit de la dependencia \pkcode{yield}. & Ese commit ocurre \emph{después} de responder, así que un refetch inmediato puede leer la base antes que él y devolver el estado viejo. Con \pkcode{staleTime: 0} en el cliente, esa carrera es la norma y no la excepción. \\
  \pkid{AD-05} & El chat de la aplicación es un cliente MCP más. & Que el endpoint de chat llame a \pkpath{services/} directamente. & Una sola superficie de datos y un solo modelo de permisos. El agente no puede hacer nada que un asistente externo no pueda, y actúa con el token del usuario, no como servicio. \\
  \pkid{AD-06} & Aceptar un intercambio registra el acuerdo pero no mueve cartas. & Transferir filas de \pkcode{collection\_item}. & La colección es lo que alguien tiene físicamente; sólo esa persona sabe si el intercambio ocurrió. Mover filas convertiría un acuerdo en una mentira sobre el mundo. \\
  \pkid{AD-07} & Un sobrante se calcula (\pkcode{SUM(quantity) > 1}), no se marca. & Una columna «disponible para cambio». & Un estado marcado se desincroniza de las cantidades reales. Calculado, el tablón nunca puede prometer una carta que ya no existe, y se revalida al publicar, al ofertar y al aceptar. \\
  \pkid{AD-08} & \pkcode{card\_price} guarda una fila por carta y por día. & Guardar sólo el último precio en \pkcode{card}. & «Un portafolio sin serie es un número suelto»: todo panel de variación reportaba +0.0\% sobre un único punto. El upsert sobre \pkcode{(card\_id, recorded\_on)} mantiene un punto por día y hace el trabajo idempotente. \\
  \pkid{AD-09} & La lectura de precios es una tarea suelta al arrancar, no un planificador. & Cron, o un worker aparte. & «Un catálogo inalcanzable es una serie desactualizada, no una aplicación que no arranca.» La deduplicación por día de calendario hace irrelevante cuántas veces se reinicie el servicio. \\
  \pkid{AD-10} & Una publicación es una oferta sin destinatario; tomarla escribe una oferta ya aceptada. & Un modelo de subasta o de reserva. & Tomarla \emph{es} el acuerdo. La reserva atómica (\pkcode{UPDATE … WHERE status='open' … RETURNING}) resuelve la concurrencia sin bloqueos, y desde ahí un trato del tablón es idéntico a uno negociado. \\
  \pkid{AD-11} & La condición se copia sobre la carta de la oferta. & Referenciar la fila de colección de origen. & Una oferta es una promesa sobre una carta en un estado, y debe sobrevivir a que su dueño edite o borre esa fila. \\
  \pkid{AD-12} & Un hilo directo es un par ordenado con \pkcode{CHECK}. & Emisor y destinatario, o una tabla de participantes. & Con un par no ordenado, dos personas escribiéndose a la vez abren dos hilos y cada lectura tiene que unirlos. No hay tabla de participantes porque no hay un tercero. \\
  \pkid{AD-13} & Los filtros escriben la URL con la History API. & \pkcode{router.replace}. & En Next eso es una navegación aunque sólo cambie un parámetro: refetch del payload de ruta y remontado del árbol. Cada clic en un filtro parpadeaba con un esqueleto y volvía a pedir lo que ya estaba en pantalla. El precio —que \pkcode{useSearchParams} deja de enterarse— se paga concentrando la lectura en un solo hook. \\
  \pkid{AD-14} & El cliente HTTP se genera del OpenAPI y se versiona en git. & Escribirlo a mano, o generarlo en tiempo de construcción. & El contrato lo fija el servidor; el frontend no puede desviarse en silencio. Versionarlo hace que un cambio de API aparezca como diff revisable. Su envoltorio se inyecta por llamada porque es lo que Kubb permite. \\
  \pkid{AD-15} & Las pruebas de la API corren contra Postgres real con rollback por prueba. & SQLite en memoria, o dobles del repositorio. & La mitad de las invariantes viven en la base: enums, \pkcode{CHECK}, \pkcode{ON CONFLICT … NULLS NOT DISTINCT}, \pkcode{RETURNING}. Un doble las daría todas por buenas. Y son rápidas igualmente: 292 casos en 13.6 s. \\
  \pkid{AD-16} & Las pruebas de la web simulan la red en la frontera de \pkcode{fetch}. & Sustituir el cliente generado. & Sustituirlo dejaría sin probar el único código propio de esa capa: el orden de cabeceras, la caché del token y la traducción de errores de FastAPI. \\
  \pkid{AD-17} & El \emph{toolset} MCP del agente está escrito a mano. & Usar el cliente MCP de pydantic-ai. & Ese cliente apunta al SDK 1.x y ni siquiera importa contra 2.0, que renombró los campos del protocolo y movió \pkcode{RequestContext}. \\
  \pkid{AD-18} & Las imágenes de escaneo van a un volumen local tras un \pkcode{Protocol}. & MinIO o S3. & «Añadiría un cuarto contenedor a algo que sólo corre en local.» La costura está puesta: cambiar de backend es implementar tres métodos. \\
  \pkid{AD-19} & \pkcode{operationId} es el nombre de la función manejadora. & El identificador por omisión de FastAPI. & El de fábrica produce ganchos como \pkcode{useListCollectionApiV1CollectionGet}. \\
  \pkid{AD-20} & Un solo tema, claro, forzado. & Tema oscuro, o seguir al sistema operativo. & Se retiró el tema oscuro a petición del usuario, y \emph{forzarlo} en vez de poner un valor por omisión: un valor por omisión sigue obedeciendo al sistema operativo y dejaría a una máquina en oscuro dentro de un tema que la aplicación ya no tiene. Costó un bloque de variables porque la paleta estaba escrita como tokens semánticos —superficies y bordes, nunca un gris literal. \\
\end{pktable}

\begin{pknote}[Una nota de mantenimiento]
El bloque \pkcode{:root} de \pkpath{packages/ui/src/styles/globals.css} conserva
un comentario de la etapa anterior («dark is the designed state») que los
valores de debajo ya contradicen. No afecta a nada en ejecución, pero conviene
saberlo antes de leer ese archivo buscando la intención de diseño: la intención
vigente es la de \pkid{AD-20}.
\end{pknote}
